% Auto-split to keep each TeX source < 600 lines.

\paragraph{\texorpdfstring{$E_8$}{E8} 平方谱塌缩与四阶常数的整系数指纹}\label{para:pom-s4-e8-square-spectrum-collapse}

\begin{proposition}[\texorpdfstring{$E_8$}{E8} 平方谱塌缩的四次证书、判别式与迹不变量]\label{prop:pom-e8-square-spectrum-collapse-trace7}
令
\[
\mathcal R_{E_8}:=\left\{4\cos^2\!\left(\frac{\pi k}{30}\right):\ k\in\{1,7,11,13\}\right\}.
\]
则 \(\mathcal R_{E_8}\) 的四个元素恰为同一整系数四次多项式的全部实根；其闭式证书、判别式以及由 Vieta 公式导出的初等对称和见式 \eqref{eq:collision_kernel_E8_square_spectrum_collapse}。
特别地，
\[
\sum_{r\in\mathcal R_{E_8}}r=7,
\]
从而整数 \(7\) 作为 \(\mathcal R_{E_8}\) 的迹不变量（trace invariant）以整系数方式出现。
\end{proposition}
\begin{proof}
由 \(\mathcal R_{E_8}\) 的定义可把 \(4\cos^2(\pi/30)=2+\zeta_{30}+\zeta_{30}^{-1}\) 视为 \(\QQ(\zeta_{30})^+\) 中的代数整数。
对其求最小多项式并计算判别式即可得到式 \eqref{eq:collision_kernel_E8_square_spectrum_collapse} 的四次证书与判别式分解；并且该四次证书表明其生成的数域次数为 \(4\)，与 \(\QQ(\zeta_{30})^+\) 的次数一致。
对称和由该四次多项式的系数通过 Vieta 公式直接读出。证毕。
\end{proof}

\begin{proposition}[\texorpdfstring{$E_8$}{E8} 指数交点的域判别式极小性与 \texorpdfstring{$t_{E_8}$}{tE8} 的指数分解]\label{prop:pom-e8-intersection-field-discriminant-index}
令 \(r_{E_8}:=4\cos^2(\pi/30)\)，并置 \(K_{E_8}:=\QQ(r_{E_8})\)。
则 \(K_{E_8}\) 为四次全实域，且其域判别式满足
\[
\mathrm{Disc}(K_{E_8})=1125=3^2\cdot 5^3
\]
（见式 \eqref{eq:collision_kernel_E8_square_spectrum_collapse} 中的 cyclotomic 证书），并且
\[
\mathcal O_{K_{E_8}}=\ZZ[r_{E_8}].
\]
另一方面，令 \(t_{E_8}=t(I_{E_8})\) 为四阶核族 \(A_4(t)\) 的 \(E_8\) 交点参数（命题 \ref{prop:pom-a4t-jones-ade-intersections}），则 \(t_{E_8}\in K_{E_8}\) 且 \(\QQ(t_{E_8})=K_{E_8}\)，但其幂整子环一般不为极大整环；
并且指数满足闭式
\[
[\mathcal O_{K_{E_8}}:\ZZ[t_{E_8}]]=5^2\cdot 29\cdot 59\cdot 1229.
\]
因此
\[
\mathrm{Disc}(\ZZ[t_{E_8}])=\mathrm{Disc}(K_{E_8})\cdot[\mathcal O_{K_{E_8}}:\ZZ[t_{E_8}]]^{2},
\]
从而与 \(t_{E_8}\) 相关的“巨大判别式”除 \(3,5\) 外的素因子并非分歧，而来自幂整生成元的指数伪影。
\end{proposition}
\begin{proof}
由 \(r_{E_8}=4\cos^2(\pi/30)=2+\zeta_{30}+\zeta_{30}^{-1}\) 可知 \(K_{E_8}\subseteq \QQ(\zeta_{30})^+\)；
并且由式 \eqref{eq:collision_kernel_E8_square_spectrum_collapse} 的四次证书可知 \([K_{E_8}:\QQ]=4\)，而 \([\QQ(\zeta_{30})^+:\QQ]=\tfrac12[\QQ(\zeta_{30}):\QQ]=4\)，故 \(K_{E_8}=\QQ(\zeta_{30})^+\)。
式 \eqref{eq:collision_kernel_E8_square_spectrum_collapse} 同时给出 \(\mathrm{Disc}(\QQ(\zeta_{30}))\) 及其平方根 \(\mathrm{Disc}(\QQ(\zeta_{30})^+)\)，并与 \(r_{E_8}\) 的四次证书判别式一致，从而 \(\ZZ[r_{E_8}]\) 的指数为 \(1\)，即为极大整环。
\par
另一方面，\(t_{E_8}=t(r_{E_8})\) 为 \(r_{E_8}\) 的有理函数，故 \(t_{E_8}\in K_{E_8}\)；而式 \eqref{eq:collision_kernel_E8_square_spectrum_collapse} 给出 \(t_{E_8}\) 的整系数最小多项式 \(P_{E_8}(t)\)，从而 \(\ZZ[t_{E_8}]\) 为 \(K_{E_8}\) 的一幂整子环，且
\(\mathrm{Disc}(\ZZ[t_{E_8}])=\mathrm{Disc}_t(P_{E_8})\)。
由判别式—指数恒等式
\(
\mathrm{Disc}(\ZZ[\alpha])=\mathrm{Disc}(K)\,[\mathcal O_K:\ZZ[\alpha]]^2
\)
取 \(\alpha=t_{E_8}\) 即得指数闭式以及“巨大判别式”的分解。证毕。
\end{proof}

\begin{corollary}[\texorpdfstring{$q=4$}{q=4} 的 cyclotomic--ADE 层与 Newman--临界层在二次分辨率上的实/虚分离]\label{cor:pom-q4-arithmetic-two-layer-real-imag}
在四阶核族 \(A_4(t)\) 的 \(E_8\) 交点处，由命题 \ref{prop:pom-e8-intersection-field-discriminant-index} 得到的 cyclotomic--ADE 层为全实域；
而 Newman 临界层的判别式二次子域为虚二次域 \(\QQ(\sqrt{-1151})\)（命题 \ref{prop:pom-a4t-newman-k2-class-number}）。
因此两层在二次分辨率上必然无交：其二次子域的交仅为 \(\QQ\)。
从而在模 \(p\) 的局部审计中，可以用两条相互正交的探针分离两类机制：一条以 cyclotomic--ADE 层的小模剩余类结构为主，另一条以临界判别式的二次特征（Kronecker 符号）为主。
\end{corollary}
\begin{proof}
全实域不含任何非平凡虚二次子域。
由于 \(\QQ(\sqrt{-1151})\) 为虚二次域，故其不可能嵌入 cyclotomic--ADE 层；
从而两者在二次分辨率层面交仅为 \(\QQ\)。
最后一句为该分离事实在有限域指纹设计中的直接翻译：实层探针由同余类/扩域次数控制，虚层探针由 \(\left(\frac{\mathrm{Disc}}{p}\right)\) 之类的二次特征控制。证毕。
\end{proof}

\begin{proposition}[四阶核族的谱自对偶冻结与四条 \texorpdfstring{$t$}{t}-不变量]\label{prop:pom-a4t-spectral-selfduality-invariants}
令四阶五态核族的特征多项式写为
\[
P_t(\lambda)=\lambda^5-2\lambda^4-t\lambda^3-2\lambda+2.
\]
则存在与 \(t\) 无关的严格恒等式
\[
\boxed{\;P_t(\lambda)+P_t(-\lambda)=4(1-\lambda^4)\;}
\]
等价地，\(P_t\) 的偶部完全冻结，而 \(t\) 仅改变其奇部的单一系数。
从而若 \(\{\lambda_i(t)\}_{i=1}^5\) 为全部特征值，则恒有四条 \(t\)-不变量：
\[
\sum_{i=1}^5 \lambda_i(t)=2,\qquad
\prod_{i=1}^5 \lambda_i(t)=-2,\qquad
\sum_{i=1}^5 \frac{1}{\lambda_i(t)}=1,\qquad
\sum_{1\le i<j\le 5}\frac{1}{\lambda_i(t)\lambda_j(t)}=0.
\]
因此四阶 Newman 相变并非由高维参数调谐触发，而发生在一个被自对偶约束锁定的谱截面上，仅凭单一自由度 \(t\) 即可完成越线。
\end{proposition}
\begin{proof}
将 \(P_t(-\lambda)\) 展开并与 \(P_t(\lambda)\) 相加，\(\lambda^5\) 与 \(t\lambda^3\) 项相消，得到
\(
P_t(\lambda)+P_t(-\lambda)=4-4\lambda^4=4(1-\lambda^4)
\)。
其余结论由 Vieta 公式给出：由于 \(P_t(\lambda)=\prod_i(\lambda-\lambda_i(t))\)，故
\(\sum_i\lambda_i(t)=2\) 与 \(\prod_i\lambda_i(t)=-2\) 由 \(\lambda^4\) 与常数项读出；
同时
\(
\sum_i \lambda_i(t)^{-1}=\frac{\sum_{i<j<k<\ell}\lambda_i\lambda_j\lambda_k\lambda_\ell}{\prod_i\lambda_i}
=\frac{-2}{-2}=1
\)，
以及
\(
\sum_{i<j}(\lambda_i\lambda_j)^{-1}
=\frac{\sum_{i<j<k}\lambda_i\lambda_j\lambda_k}{\prod_i\lambda_i}
=0
\)，
由 \(\lambda\) 与 \(\lambda^2\) 系数读出。证毕。
\end{proof}

\input{sections/generated/eq_collision_kernel_E8_square_spectrum_collapse}

